\section{Die Arbeit}
\newvideofile{Arbeit}{Arbeit}
	\s{
	In diesem Kapitel wird zur Unterscheidung zwischen Arbeit und Energie das Symbol $E$ für Energie (ohne Vektorpfeil) und $W$ für Arbeit (vom englischen Wort "work") verwendet.
	Obwohl das Formelzeichen $W$ umgangssprachlich häufig sowohl für Arbeit als auch für Energie benutzt wird, ist Arbeit im Wesentlichen eine Veränderung der Energie. Beide Größen,
	Arbeit und Energie, werden in der Einheit Joule (J) gemessen.
	Wie im Abschnitt über die Energieerhaltung beschrieben, besagt der Energieerhaltungssatz, dass Energie weder erzeugt noch vernichtet werden kann.
	Sie ändert lediglich ihre Form. % Dieser Prozess der Umwandlung wird auch als Arbeit bezeichnet.
	Diesen Zusammenhang beschreibt die folgende Formel:


	\begin{equation}
		W = \Delta E \quad \quad [W] = \text{1 Joule = 1 J = 1 Nm}
	\end{equation}}

\begin{frame}[t]
	\ftx{Die Arbeit}\index{Arbeit}
\b{

	\begin{itemize}
		\item Arbeit ist die Änderung der Energieform
	\end{itemize}

	\begin{equation*}
		W = \Delta E \quad \quad [W] = \text{1 Joule = 1 J = 1 Nm}
	\end{equation*}
	\vspace*{1cm}
	\only<2->{
			\begin{itemize}
				\centering \item $ \displaystyle W = - \int_{P_\mathrm{1}}^{P_\mathrm{2}} \vec{F} \cdot \mathrm{d}\vec{s} $
			\end{itemize}}
}% nur Folien

\speech{02_Arbeit1}{1}{
	Wir haben jetzt verschiedene Arten der Energie kennengelernt. Wir haben gesehen, wie wir Energie berechnen können, nämlich über das Linienintegral. Und letztlich hatten wir bei der Einführung am Beispiel des Pumpspeicherkraftwerks auch den Begriff der Arbeit definiert. Arbeit ist nichts anderes als die Änderung einer Energieform. Das heißt, die Arbeit, die wir verrichten, ist das Delta, also die Differenz der Energie vorher und nachher. Die Einheit für die Arbeit ist wie bei der Energie auch Joule, beziehungsweise Newtonmeter. Das heißt, wenn wir jetzt nochmal an das Pumpspeicherkraftwerk von gerade denken, dann haben wir eine gewisse potenzielle Energie im Oberbecken gespeichert. Diese ist abhängig von der Höhe zwischen Ober- und Unterbecken, und der Menge an Wasser, dem Volumen, beziehungsweise damit verknüpft, dem Gewicht. Das ist die Energie. Und wenn wir dieses Wasser durch die Turbine laufen lassen, dann wird Arbeit verrichtet. Und zwar genau so viel Arbeit, wie wir Energie aus dem Oberbecken in das Unterbecken verlagern. Das ist die Arbeit, die wir verrichten.
}

\s{
	Diese Änderung der Energiemenge kann durch verschiedene physikalische Aktionen hervorgerufen werden, einschließlich der Verschiebung eines Objekts in einem Kraftfeld.
	Um diese Konzepte weiter zu veranschaulichen, ist im Folgenden die allgemeine Formel \ref{eq:WFds} für die Arbeit dargestellt:


	\begin{equation}
		W = - \int_{P_\mathrm{1}}^{P_\mathrm{2}} \vec{F} \cdot \mathrm{d}\vec{s}
		 \label{eq:WFds}
	\end{equation}}


	%%% Begründung des negativen Vorzeichens %%%
\s{
Die Gleichung berechnet Arbeit durch Integration der Kraft entlang eines Weges, wobei die Kraftausrichtung zur
Bewegungsrichtung des Objekts entscheidend ist. Ein Vorzeichenwechsel in der Formel deutet darauf hin, dass Energie aus
einem Energiespeicher für Arbeit verwendet wird. Dieses Vorzeichen ist kontextabhängig und ändert sich je nach Perspektive.
Wird Energie entnommen, gilt die Arbeit als negativ, da sie von der gespeicherten Gesamtenergie abgezogen wird,
was eine Verringerung der Systemenergie signalisiert.
}
\b{\vspace*{-0.5cm}}
\only<3->{
	\begin{Merksatz}{Arbeit}
	 Das Vorzeichen der Arbeit ist kontextabhängig: % und beschreibt, ob die Arbeit ein Energieniveau anhebt oder absenkt.
	Die Arbeit ist positiv, wenn Energie einem System hinzugefügt wird
	und negativ, wenn Energie einem System entnommen wird.

	\end{Merksatz}}
\speech{02_Arbeit2}{2}{
Demzufolge können wir das formelmäßig genauso beschreiben. Die Arbeit W ist gleich dem Linienintegral von P1 nach P2 über F d s.
}
\speech{02_Arbeit2}{3}{Merke dir also dass das Vorzeichen der Arbeit abhängig ist ob einem System Energie hinzugefügt wird dann ist es Positiv, oder ob Sie entnommen wird dann ist es Negativ.
}
\end{frame}


	\subsection{Die elektrische Arbeit}
\s{
	In der Elektrostatik sind die Kräfte $\vec{F}_{\mathrm{el}}$, die zwischen Ladungsträgern $Q$ und elektrischen Feldern $\vec{E}$ wirken, zentral.
	Diese Kräfte entstehen, wenn elektrische Felder auf Ladungsträger wirken.
	Das elektrische Feld $\vec{E}$ wird als Kraft pro Ladung definiert und beschreibt,
	wie groß die Kraft $\vec{F}_{\mathrm{el}}$ im Verhältnis zur eingebrachten Ladung $Q$ ist.

}
\begin{frame}
	\ftx{Die elektrische Arbeit}
	\begin{overlayarea}{\textwidth}{0.9\textheight}

	\b{\only<1->{
		\begin{itemize}\centering
  			\item $\displaystyle \vec{E} = \frac{\vec{F}_\mathrm{el}}{Q}
  			\quad \Rightarrow \quad
  			\vec{F}_\mathrm{el} = Q \cdot \vec{E}$
		\end{itemize}

		\begin{itemize}\centering
  			\item $\displaystyle W_\mathrm{el} = \int_{P_\mathrm{1}}^{P_\mathrm{2}} - Q \cdot \vec{E} \cdot \mathrm{d}\vec{s}
			\quad \Rightarrow \quad
			 W_\mathrm{el} = - Q \int_{P_\mathrm{1}}^{P_\mathrm{2}} \vec{E} \cdot \mathrm{d}\vec{s} $
		\end{itemize}

		\begin{itemize}\centering
  			\item $\displaystyle W_\mathrm{{el}} = Q \int_{P_\mathrm{1}}^{P_\mathrm{2}} \vec{E} \cdot \mathrm{d}\vec{s}
			\quad \Rightarrow \quad
			W_\mathrm{{el}} = Q \cdot U $
		\end{itemize}

	\begin{itemize}\centering
  			\item $\displaystyle I = \frac{\mathrm{d}Q}{\mathrm{d}t}
			 \quad \Rightarrow \quad
			  	Q = I \cdot t $
		\end{itemize}}

}

\speech{02_Arbeit3}{1}{
Jetzt betrachten wir die elektrische Arbeit. Wie groß ist die Arbeit, die wir mit elektrischen Systemen verrichten? Wir fangen genauso an wie vorher auch. Wir betrachten die elektrische Energie: E ist gleich F-EL durch Q, wodurch folgt F-EL ist gleich Q mal E. Daraus abgeleitet haben wir dann die elektrische Arbeit, die letztlich minus Q mal dem Integral von P1 nach P2 über E mal D-S entspricht.
Wenn wir jetzt die elektrische Arbeit, die wir verrichten, mit F-EL gleich Q mal E gleichsetzen, dann erhalten wir durch Einsetzen oder durch Vergleichen mit Formeln, die wir bereits vorher hergeleitet haben, dass das Integral von P1 nach P2 über E mal D-S nichts anderes ist als unsere definierte Spannung U. Das heißt, die elektrische Arbeit W-EL, die wir verrichten, ist Q mal U, die Ladung mal die Spannung.
Wenn wir dann weiter in die Definition des Stroms schauen, dann hatten wir ja gesagt, dass der Strom I die zeitliche Veränderung der Ladungsmenge Q ist. I ist also gleich d-Q durch d-t. Es folgt also letztlich, dass wir die Ladung Q beschreiben können als den Strom I mal die Zeit T, die der Strom fließt.
}


\s{
	\begin{equation}
		\vec{E} = \frac{\vec{F}_\mathrm{el}}{Q} \quad \Rightarrow \quad \vec{F}_\mathrm{el} = Q \cdot \vec{E}
	\end{equation}

	Wenn nun die Formel für die elektrische Kraft $\vec{F}_\mathrm{el}$ in die Formel \ref{eq:WFds} für die allgemeine Arbeit eingesetzt wird, ergibt das die elektrische Arbeit $W_\mathrm{el}$.
	Da sich die Ladung $Q$ über die Strecke $s$ nicht verändert, kann diese vor das Integral gezogen werden.
	Da das elektrische Feld $\vec{E}$ ein Vektor ist und somit einen Betrag und eine Richtung hat, ist das Produkt aus $\vec{E}$ und $\mathrm{d}\vec{s}$ in jedem Fall von der gewählten Strecke $s$ abhängig.
	Daher muss über das Produkt $\vec{E}$ und $\mathrm{d}\vec{s}$ ein Linienintegral gebildet werden.

	\begin{equation}
		W_\mathrm{el} = \int_{P_\mathrm{1}}^{P_\mathrm{2}} - Q \cdot \vec{E} \cdot \mathrm{d}\vec{s}  \quad \Rightarrow \quad W_\mathrm{el} = - Q \int_{P_\mathrm{1}}^{P_\mathrm{2}} \vec{E} \cdot \mathrm{d}\vec{s}
	\end{equation}

	Aus Modul 1 ist bekannt, dass das Integral über $\vec{E} \cdot \mathrm{d}\vec{s}$ im konstanten Fall die Spannung $U$ ergibt.
	Für eine verbraucherunabhängige Darstellung wird in der folgenden Formel das negative Vorzeichen weggelassen, sowie für den konstanten Fall durch Einsetzen der Spannung $U$ vereinfacht dargestellt.

	\begin{equation}
		W_\mathrm{{el}} = Q \int_{P_\mathrm{1}}^{P_\mathrm{2}} \vec{E} \cdot \mathrm{d}\vec{s} \quad \Rightarrow \quad W_\mathrm{{el}} = Q \cdot U\label{eq:WQU}
		\end{equation}


	Da für einen elektrischen Energietransport eine elektrische Stromstärke notwendig ist, kann nun der aus Modul 1 bekannte Zusammenhang zwischen elektrischer Stromstärke $I$ und der Ladung $Q$ verwendet und umgestellt werden,
	um die Ladung $Q$ in der Formel \ref{eq:WQU} auszutauschen. Ein konstanter Stromfluss $I$ ergibt sich aus einer gleichmäßigen Bewegung der Ladungsträger $Q$ über die Zeit $t$.
	In diesem Fall kann auf den Differentialoperator ($\mathrm{d}$) verzichtet werden. Durch Umstellung nach $Q$ ergibt sich folgende Formel:
	\begin{equation}
		I = \frac{\mathrm{d}Q}{\mathrm{d}t} \quad \Rightarrow \quad 	Q = I \cdot t
	\end{equation}

	Eingesetzt ergibt dies nun folgenden Ausdruck, welcher verdeutlicht, dass das Verrichten oder Aufbringen von Arbeit einen Stromfluss voraussetzt.

	\begin{equation}
		W_\mathrm{{el}} = U \cdot I \cdot t \quad \quad [W_\mathrm{el}] = \text{1 Joule = 1 J = 1 Ws}\label{eq:WUIt}
	\end{equation}
}
\b{
\only<2->{\begin{itemize}
	\item Eingesetzt ergibt dies nun folgenden Ausdruck
		  \end{itemize}

	\begin{equation*}
		W_\mathrm{{el}} = U \cdot I \cdot t \quad \quad [W_\mathrm{el}] = \text{1 Joule = 1 J = 1 Ws}
	\end{equation*}}

		\only<3->{ \vspace{-7.9cm}
		\begin{tikzpicture}{remember picture,overlay}
		\fill[white,opacity=0.7](current page.south west) rectangle(current page.north east);
		\end{tikzpicture}
		\vspace{-8.6cm}
		\begin{Merksatz}{Arbeit}

			\begin{itemize}
				\item Energie beschreibt das Potential Arbeit zu verrichten.
				\item Energie kann gespeichert werden, Arbeit nicht.
				\item Für das Verrichten einer Arbeit wird immer eine Zeitdauer benötigt.
				\item Ein Momentanzustand ist immer der Ist-Zustand der Energieverteilung.
				\item Das Verrichten elektrischer Arbeit bedarf immer eines Stromflusses.
			\end{itemize}

		\end{Merksatz}}}
	\end{overlayarea}
		\s{\begin{Merksatz}{Arbeit}

			\begin{itemize}
				\item Energie beschreibt das Potential Arbeit zu verrichten.
				\item Energie kann gespeichert werden, Arbeit nicht.
				\item Für das Verrichten einer Arbeit wird immer eine Zeitdauer benötigt.
				\item Ein Momentanzustand ist immer der Ist-Zustand der Energieverteilung.
				\item Das Verrichten elektrischer Arbeit bedarf immer eines Stromflusses.
			\end{itemize}

		\end{Merksatz}}


\speech{02_Arbeit3}{2}{ %seed imer noch nicht perfekt, aber schon ok
Wir ersetzen also das Q in der Hergeleiteten Formel durch I mal T und erhalten: W-EL ist gleich U mal I mal T. Auch hier ist die Einheit wieder Joule oder auch Wattsekunde.}

\speech{02_Arbeit3}{3}{Merke also das Energie das Potential Arbeit zu verrichten beschreibt. Energie kann gespeichert werden Arbeit nicht. Für das Verrichten von Arbeit wird immer eine Zeitdauer benötigt. Der Momentanzustand ist immer der Ist-Zustand der Energieverteilung, wo also, wie viel Energie im System gespeichert ist. Und wenn wir elektrische Arbeit verrichten wollen, dann brauchen wir immer einen Stromfluss. Wenn wir keinen Stromfluss haben, wird keine Arbeit verrichtet. Genauso übrigens, wenn keine Spannung anliegt.
}
\end{frame}



\subsection{Wegintegral der elektrischen Arbeit}
\s{
	Die Bewegung einer Ladung im elektrischen Feld ist lediglich abhängig von der Positionsänderung entlang der elektrischen Feldlinien.
	Abhängig von der Bewegungsrichtung entlang oder entgegen der Feldlinien ergibt das entsprechende Vorzeichen.
	Bewegungen entlang einer Äquipotentialfläche (Fläche mit identischem elektrischen Potential) haben keinen Einfluss auf die elektrische Arbeit.
	In Abbildung \ref{fig:Ladung_E-Feld_Arbeit} sind die Punkte $P_\mathrm{1}$ und $P_\mathrm{2}$ auf je unterschiedlichen Potentialen dargestellt, welche jeweils durch zwei unterschiedliche Strecken $s_\mathrm{1}$ und $s_\mathrm{2}$ verbunden sind.
}
\begin{frame}
	\ftx{Wegintegral der elektrischen Arbeit}
	\begin{overlayarea}{\textwidth}{0.9\textheight}
	\s{
	\begin{figure}[H]

		\includesvg[width=1\textwidth]{Ladung_E-Feld_Arbeit}
		\caption{\textbf{Bewegung einer Punktladung.} Zeigt die Bewegung einer Punktladung entlang zweier Strecken. Die Arbeit die verrichtet wird entlang einer geschlossenen
		 Strecke in einem statischen elektrischen Feld ist Null.}
		 \alttext{Die Abbildung zeigt eine Punktladung Q, die sich in einem statischen elektrischen Feld bewegt. Die roten horizontalen Pfeile symbolisieren das elektrische Feld, das von links nach rechts gerichtet ist. Zwei Wege zwischen Punkt P eins und Punkt P zwei sind eingezeichnet. Ein Weg verläuft relativ direkt zwischen beiden Punkten, der andere beschreibt eine längere, gekrümmte Bahn. Die Darstellung macht deutlich, dass in einem statischen elektrischen Feld die verrichtete Arbeit entlang einer geschlossenen Strecke gleich Null ist, unabhängig davon, welche konkrete Bahn gewählt wird.}
		\label{fig:Ladung_E-Feld_Arbeit}
	\end{figure}}

	\b{\begin{figure}[H]
		\includesvg[width=1\textwidth]{Ladung_E-Feld_Arbeit}
	\end{figure}}


	\s{
	Die Beträge zweier äquidistanter Positionsänderungen in entgegengesetzen Richtungen entlang der elektrischen Feldlinien sind identisch.
	Diese Eigenschaften führen dazu, dass die Positionsänderung einer Ladung $Q$ von einem Punkt $P_\mathrm{1}$ zu einem Punkt $P_\mathrm{2}$ betragsmäßig gleich einer beliebig anderen Strecke von $P_\mathrm{2}$ zu $P_\mathrm{1}$ entspricht.


	\begin{equation}
		\left| W_\mathrm{el} \right| = \left| -Q \cdot \int_{s_\mathrm{1}} \vec{E} \cdot \mathrm{d}\vec{s}\; \right| = \left| -Q \cdot \int_{s_\mathrm{2}}\vec{E} \cdot \mathrm{d}\vec{s}\; \right| %=
	\end{equation}

		Dieses Verhalten spiegelt wieder, dass die Arbeit gleich groß ist, jedoch in einem Fall Arbeit aufgewendet und im anderen Fall Arbeit verrichtet wird.
		Dies führt zu entgegengesetzten Vorzeichen, was dazu führt, dass sich die Summe beider Arbeiten aufhebt.
		Mathematisch kann dies auch so ausgedrückt werden, dass das Wegintegral entlang einer beliebigen Kontur zurück zum Ausgangspunkt Null ergibt.
		Da sich die Ladung $Q$ außerhalb des Umlaufintegral befindet, ist das Ergebnis somit unabhängig von der Ladungsmenge.

	\begin{equation}
	W_\mathrm{el} = - Q \cdot \int_{s_\mathrm{1}} \vec{E} \cdot \mathrm{d}\vec{s} -Q \cdot \int_{s_\mathrm{2}}\vec{E} \cdot \mathrm{d}\vec{s} = - Q \cdot \oint_{s}\vec{E} \cdot \mathrm{d}\vec{s} = 0
	\rightarrow \oint_{s}\vec{E} \cdot \mathrm{d}\vec{s} = 0
	\end{equation}

Das Ergebnis des Wegintegrals verdeutlicht ein fundamentales Prinzip der Physik: In einem statischen Feld ist die über einen geschlossenen Weg verrichtete Arbeit immer Null. Dies bestätigt die Wegunabhängigkeit der Arbeit in diesen Feldern und zeigt, dass die verrichtete Arbeit ausschließlich von den Anfangs- und Endpunkten abhängt und nicht vom spezifischen Weg zwischen diesen Punkten.

\begin{Merksatz}{Wegintegral}
Das geschlossene Wegintegral $ \displaystyle \oint_{s}\vec{E} \cdot \mathrm{d}\vec{s} $ ist gleich Null.
\end{Merksatz}


}% nur skript
\b{
	\only<2->{
	\begin{itemize} \centering
		\item $\displaystyle \left| W_\mathrm{el} \right| = \left| -Q \cdot \int_{s_\mathrm{1}} \vec{E} \cdot \mathrm{d}\vec{s}\; \right| = \left| -Q \cdot \int_{s_\mathrm{2}}\vec{E} \cdot \mathrm{d}\vec{s}\; \right| $
		\item  $\displaystyle W_\mathrm{el} = - Q \cdot \int_{s_\mathrm{1}} \vec{E} \cdot \mathrm{d}\vec{s} -Q \cdot \int_{s_\mathrm{2}}\vec{E} \cdot \mathrm{d}\vec{s} = - Q \cdot \oint_{s}\vec{E} \cdot \mathrm{d}\vec{s} = 0
	\rightarrow \oint_{s}\vec{E} \cdot \mathrm{d}\vec{s} = 0 $
	\end{itemize}}

	\only<3->{
	\vspace{-7.9cm}
		\begin{tikzpicture}{remember picture,overlay}
		\fill[white,opacity=0.7](current page.south west) rectangle(current page.north east);
		\end{tikzpicture}
		\vspace{-8.5cm}
	\begin{Merksatz}{Wegintegral}

	Das geschlossene Wegintegral $ \displaystyle \oint_{s}\vec{E} \cdot \mathrm{d}\vec{s} $ ist gleich Null.
	\end{Merksatz}}

}
\end{overlayarea}
	\speech{02_Arbeit4}{1}{
	Schauen wir uns jetzt an, wie elektrische Arbeit im elektrischen Feld beschrieben wird.
	Auf eine Ladung wirkt im Feld eine Kraft. Bewegen wir die Ladung ein kleines Stück weiter, also um das Wegstück ds, dann ist die Arbeit Q mal das Integral über das Skalarprodukt von E mal ds. Weil Kraft und Weg Richtungen haben, also Vektoren sind, zählt für die Arbeit nur der Anteil der Kraft in Bewegungsrichtung, deshalb verwendet man das Skalarprodukt.
	Man zerlegt den Weg in viele kleine Stücke ds und summiert über den ganzen Weg auf. Das ist das Linien- oder Wegintegral.
	Jetzt zur Kernaussage der Folie.}
	\speech{02_Arbeit4}{2}{
	Wir können von P1 nach P2 über den Weg S1 oder über einen anderen Weg S2 gehen.
	Im elektrostatischen Feld ist die verrichtete Arbeit gleich, sie hängt nur von Start- und Endpunkt ab, nicht vom Weg S1 oder S2. Das zeigt die obere Gleichung.
	Wenn wir den Weg schließen, also von P1 nach P2 über S1 und von P2 zurück nach P1 über S2, entsteht ein geschlossenes Wegintegral. Das erkennt man am Integralzeichen mit dem Kreis.
	Für ein elektrostatisches Feld gilt dann: Die Gesamtarbeit über den geschlossenen Weg ist null. Deshalb ist das Wegintegral von E mal ds gleich null.
	Das Minuszeichen vor Q hängt mit der Vorzeichenkonvention zusammen: Je nachdem, ob wir die Arbeit des Feldes betrachten oder die Arbeit, die wir von außen aufbringen, ändert sich das Vorzeichen. Und wichtig als Einschränkung: Diese Wegeunabhängigkeit gilt nur für elektrostatische, zeitlich konstante Felder.
	}
	\speech{02_Arbeit4}{3}{Wir sehen also: Im elektrostatischen Feld ist das geschlossene Wegintegral E mal ds gleich null.}
\end{frame}
